% ===================================================================
%  ТАБЛИЦЯ № 15 — Ринок. — Харчові припаси.
%  Traduction 2026 d'apres texto/io, controlee sur texto/fr, texto/en
%  et texto/es. Consigne : voir texto/uk/10-tablycia-01.tex.
%
%  LE TABLEAU DU MARCHE EST CELUI DE LA CUISINE POPULAIRE, et il rend
%  a l'ukrainien ce que les tableaux de la gare et du port lui avaient
%  pris : le boudin est une « кров'янка », le saindoux du « смалець »,
%  la miche une « паляниця », le fer a repasser une « праска », la
%  lentille de la « сочевиця ». Ce sont les mots de la maison, et ils
%  sont a la langue, non a l'administration.
% ===================================================================

%%K t15-apar-1 apar
\VUsaut{4.0mm}
\VUtitre{15.0pt}{60}{ТАБЛИЦЯ № 15}
\VUsaut{3.0mm}

%%K t15-tit-1 sub
\VUpk{11.4pt}{40}{Ринок. --- Харчові припаси.}
\VUsaut{3.0mm}

%%K t15-01-1 p
1. --- Більша частина торговців, які постачають нам потрібні припаси, зібрана
під \VUgras{повіткою} \textsuperscript{(1)} \VUgras{критого ринку} або на
сусідніх вулицях.

%%K t15-02-1 p
2. --- \VUgras{Ковбасник} \textsuperscript{(2)}, який продає \VUgras{окіст}
\textsuperscript{(3)}, \VUgras{кров'янку} \textsuperscript{(4)},
\VUgras{ковбасу} \textsuperscript{(5)}, \VUgras{сардельки}
\textsuperscript{(6)} і \VUgras{сало} \textsuperscript{(7)}, стоїть поруч із
\VUgras{торговкою маслом і сиром} \textsuperscript{(8)}, чий прилавок заставлений
\VUgras{голландськими сирами} \textsuperscript{(9)} з червоною скоринкою,
\VUgras{камамберами} \textsuperscript{(10)}, великими \VUgras{кругами ґрюєру}
\textsuperscript{(11)} і свіжими \VUgras{грудками масла} \textsuperscript{(12)}.

%%K t15-03-1 p
3. --- Перед нею \VUgras{городниця} \textsuperscript{(13)} пропонує своїм
покупцям прекрасні \VUgras{помідори} \textsuperscript{(14)}, \VUgras{цвітну
капусту} \textsuperscript{(15)}, \VUgras{салату} \textsuperscript{(16)},
\VUgras{капусту} \textsuperscript{(17)}, \VUgras{гарбузи}
\textsuperscript{(19)}, \VUgras{дині} \textsuperscript{(18)} і навіть
\VUgras{гриби} \textsuperscript{(20)}. Але \VUgras{пивовар}, який спинив свій
\VUgras{розвізний віз} \textsuperscript{(21)}, навантажений \VUgras{пивними
барильцями} \textsuperscript{(22)}, просто перед її прилавком, заважає покупцям
підійти. Городник несе їй кошик \VUgras{артишоків} \textsuperscript{(23)}.

%%K t15-tit-2 sub
\VUsaut{4.0mm}
\VUpk{10.2pt}{40}{Продавати й купувати.}
\VUsaut{3.0mm}

%%K t15-04-1 p
4. --- На ринку велике пожвавлення, бо настала година \VUgras{торгу}
\textsuperscript{(24)}. \VUgras{Господині} \textsuperscript{(26)}, які приходять
сюди по покупки, спиняються дорогою перед прилавком \VUgras{торговки
потрухами} \textsuperscript{(25)}, щоб купити \VUgras{рубці} або \VUgras{телячу
голову}.

%%K t15-05-1 p
5. --- Товстий \VUgras{кухар} \textsuperscript{(27)} торгується про прекрасного
\VUgras{тунця} з \VUgras{рибницею} \textsuperscript{(28)}, яка набавляє ціну як
їй заманеться. Їй привезли багато \VUgras{вугрів, камбали, форелі} і
\VUgras{коропів}. Її \VUgras{тріска} і її \VUgras{мідії} дуже свіжі.

%%K t15-tit-3 sub
\VUsaut{4.0mm}
\VUpk{10.2pt}{40}{Дешеве і дороге.}
\VUsaut{3.0mm}

%%K t15-06-1 p
6. --- \VUgras{Торговка птицею} \textsuperscript{(29)}, що влаштувалася поруч із
нею, дуже чесна. Продаючи \VUgras{індика}, \VUgras{гуску}, \VUgras{качку} чи
просту \VUgras{курку}, вона не просить більше за справедливу ціну.

%%K t15-07-1 p
7. --- На розі площі, на першому поверсі, недавно завела торгівлю
\VUgras{білизняниця} \textsuperscript{(30)}, у якої покупці завжди напевно
знайдуть дуже добрий вибір \VUgras{скатертин}, \VUgras{серветок, фартухів} і
\VUgras{кухонних рушників}. На тому самому поверсі влаштував свою майстерню
\VUgras{ножар} \textsuperscript{(31)}. Він гострить \VUgras{ножі} торговок і
робить городникам \VUgras{садові ножі} з найтвердішої \VUgras{сталі}.

%%K t15-tit-4 sub
\VUsaut{4.0mm}
\VUpk{10.2pt}{40}{Бакалія.}
\VUsaut{3.0mm}

%%K t15-08-1 p
8. --- Під його майстернею не так давно відкрилася велика харчова крамниця. Це
вже не проста, скромна \VUgras{бакалійна крамничка} \textsuperscript{(32)}
давніх часів, де продавали тільки вжиткові товари --- \VUgras{чай}
\textsuperscript{(33)}, \VUgras{шоколад} \textsuperscript{(34)}, \VUgras{голову
цукру} \textsuperscript{(37)} і \VUgras{мило} \textsuperscript{(38)}. Тут
продають усе: \VUgras{хрустке печиво}, \VUgras{консерви}
\textsuperscript{(35)}, \VUgras{лікери} \textsuperscript{(36)}, \VUgras{ром,
коньяк, кірш, анисівку} і \VUgras{кюрасао}. Дичину дістати тут так само легко ---
\VUgras{зайця} \textsuperscript{(39)}, \VUgras{дроздів} \textsuperscript{(40)},
\VUgras{куріпок} \textsuperscript{(41)}, \VUgras{перепілок}
\textsuperscript{(42)}, \VUgras{дику качку} \textsuperscript{(43)} чи
\VUgras{фазана} \textsuperscript{(44)} --- як і тропічні плоди на кшталт
\VUgras{бананів} \textsuperscript{(46)} і \VUgras{ананасів}.

%%K t15-09-1 p
9. --- Дуже послужливий \VUgras{продавець} \textsuperscript{(45)} готовий подати
все, що попросять: \VUgras{варення} \textsuperscript{(47)}, смородинове чи
айвове, \VUgras{смалець} \textsuperscript{(48)}, \VUgras{корнішони}
\textsuperscript{(49)} чи солоні \VUgras{сардини} \textsuperscript{(50)}.

%%K t15-09-2 p
Це дуже зручно \VUgras{покупчиням} \textsuperscript{(51)}, які знаходять поруч із
найзвичайнішими товарами найрідкісніший ранній овоч і найсмачніші плоди: соковиті
\VUgras{груші} \textsuperscript{(52)}, \VUgras{сливи} \textsuperscript{(53)},
\VUgras{виноград} \textsuperscript{(54)}, \VUgras{персики}
\textsuperscript{(55)}, \VUgras{мигдаль} \textsuperscript{(56)} і
\VUgras{горіхи} \textsuperscript{(57)}.

%%K t15-tit-5 sub
\VUsaut{4.0mm}
\VUpk{10.2pt}{40}{Покупці.}
\VUsaut{3.0mm}

%%K t15-10-1 p
10. --- \VUgras{Рис} \textsuperscript{(58)} і \VUgras{сочевиця}
\textsuperscript{(59)} стоять поруч зі скринею копчених \VUgras{оселедців}
\textsuperscript{(60)}; \VUgras{картопля} \textsuperscript{(61)} і
\VUgras{квасоля} \textsuperscript{(63)} сусідять із \VUgras{яблуками}
\textsuperscript{(62)}, \VUgras{інжиром} \textsuperscript{(64)},
\VUgras{чорносливом} \textsuperscript{(65)} і \VUgras{ліщиновими горішками}
\textsuperscript{(66)}. Свіжі \VUgras{яйця} \textsuperscript{(67)} і
\VUgras{маслини} \textsuperscript{(68)} теж є в цій крамниці, тож
\VUgras{кошики} \textsuperscript{(70)} і \VUgras{сітки} \textsuperscript{(73)}
наповнюються дуже швидко.

%%K t15-11-1 p
11. --- \VUgras{Кава} \textsuperscript{(72)} цієї добірної фірми заслужено
славиться серед \VUgras{служниць} \textsuperscript{(69)} і куховарок, бо старий
\VUgras{бакалійник} \textsuperscript{(71)} смажить її сам із якнайбільшим
дбанням.

%%K t15-tit-6 sub
\VUsaut{4.0mm}
\VUpk{10.2pt}{40}{Що тут можна побачити.}
\VUsaut{3.0mm}

%%K t15-12-1 p
12. --- Не всі приходять на ринок по припаси. У мальовничому русі вулички, що
веде до повітки, бачиш найрізніший народ. Поруч із добродушним \VUgras{різницьким
робітником} \textsuperscript{(75)}, який штовхає перед собою \VUgras{тачку}
\textsuperscript{(74)}, ось два солдати у відпустці, дуже горді своїми яскравими
мундирами: \VUgras{гусар} \textsuperscript{(76)} у синьому \VUgras{доломані} і
\VUgras{ківері із султаном} та \VUgras{капрал зуавів} у широких шароварах і з
\VUgras{фескою} на голові.

%%K t15-13-1 p
13. --- \VUgras{Пекар} \textsuperscript{(80)}, що стоїть на порозі
\VUgras{пекарні} \textsuperscript{(78)}, щойно вимісив тісто на свій
\VUgras{хліб} \textsuperscript{(79)} і вийняв із печі \VUgras{рогалики}
\textsuperscript{(81)}, \VUgras{булочки} \textsuperscript{(82)} і
\VUgras{паляниці} \textsuperscript{(83)}, що лежать у вітрині.

%%K t15-14-1 p
14. --- Над пекарнею живе вправна \VUgras{швачка} \textsuperscript{(84)}, у якої
працює кілька \VUgras{вишивальниць} \textsuperscript{(85)}. Поруч із нею живе
\VUgras{пралька і прасувальниця} \textsuperscript{(86)}, яка крохмалить
\VUgras{білизну} \textsuperscript{(88)}, виправши й висушивши її, і прасує її
\VUgras{праскою} \textsuperscript{(87)}.

%%K t15-tit-7 sub
\VUsaut{4.0mm}
\VUpk{10.2pt}{40}{М'ясо, риба і городина.}
\VUsaut{3.0mm}

%%K t15-15-1 p
15. --- У \VUgras{м'ясній крамниці} \textsuperscript{(89)} на нижньому поверсі
продають усяке м'ясо --- \VUgras{баранину} \textsuperscript{(90)},
\VUgras{яловичину} \textsuperscript{(94)} чи \VUgras{телятину}
\textsuperscript{(92)} --- у вигляді \VUgras{баранячих ніг, вирізки, печені} і
\VUgras{відбивних}. \VUgras{Різників підручний} \textsuperscript{(93)} розбирає
все це м'ясо і відкладає \VUgras{мозкові кістки} \textsuperscript{(96)}. Коло
дверей м'ясної стоїть \VUgras{торговка устрицями} \textsuperscript{(97)}, яка
продає \VUgras{устриці} \textsuperscript{(98)}. Вона їх відкриває, як попросять.

%%K t15-16-1 p
16. --- Усі ці торговці платять патент або місцеве; ось чому \VUgras{городовий}
\textsuperscript{(100)} без пощади ганяє бідних \VUgras{рознощиць городини}
\textsuperscript{(99)}, які складають їм конкуренцію, розвозячи на своїх візках
\VUgras{моркву, редиску, ріпу, цибулю-порей} або \VUgras{пучки спаржі, свіжий
горошок, шпинат, щавель, крес} і \VUgras{петрушку}.

%%K t15-tit-8 sub
\VUsaut{4.0mm}
\VUpk{10.2pt}{40}{Стережися городового!}
\VUsaut{3.0mm}

%%K t15-17-1 p
17. --- Хлопець, що торгує \VUgras{часником} \textsuperscript{(101)} і
\VUgras{цибулею}, чудово знає, що й йому не вільно продавати свій товар коло
ринку. Він пускається бігти, важачи перекинути кошик \VUgras{крабів},
\VUgras{раків} і \VUgras{креветок}, що належить \VUgras{торговці}
\textsuperscript{(102)} \VUgras{лангустами} й \VUgras{омарами}.

%%K t15-18-1 p
18. --- Усі клопочуться і зчиняють великий гамір; але це не заважає ясно чути
заклик мандрівного \VUgras{скляра} \textsuperscript{(103)} і крик
\VUgras{вугляра} \textsuperscript{(104)}, який щойно лишив на хвилину свого
візка, щоб доставити \VUgras{мішок вугілля} \textsuperscript{(105)}.
\VUsaut{6.0mm}
\VUfilet{22mm}
